\section{The \textsf{equal} Algorithm}
\label{equal}

This algorithm implements a comparison over two generic sequences.
The specification of the algorithm is:

\listingname{equal.h}
\cinput{equal/equal.spec}

The implementation is:

\listingname{equal.c}
\cinput{equal/equal.impl}

\clearpage
\subsection{Correctness}

The implementation is proved correct against its specification by
simply running the \textsf{WP} plug-in: 
\fclogs{nonmutating}{equal}

The reader should notice the warning emitted by \textsf{WP}. Actually,
the plug-in is not responsible for proving the absence of runtime
errors during program execution, since other plug-ins can be used
for this, for instance \emph{Value Analysis}.

\subsection{Safety}

However, it is still possible to completely prove the program with
\textsf{WP} thanks to \textsf{RTE} plug-in. This last plug-in
generates assertions in the program wherever a runtime error might
occur. Then, the \textsf{WP} plug-in can try to discharge the
generated assertions.

This all-together method is easily performed with \texttt{wp-rte}
option of \textsf{WP}:
\fclogs{nonmutating}{equal_rte}

In this book, we will always use this technique to prove the
correctness and the safety of studied algorithm.
